%% kapa: exam=1 round=25 year=2014 period=1 form=A subject=civil_law subject_ko=민법 seq=40 number=40
%% source: Q-Net 감정평가사 기출문제 (cst00309) artlSeq=5101207 file=extracted/1008/5101207/2014년 제25회 감정평가사 1차시험 기출문제/1교시(A형).hwp
%% conversion: kordoc 4.24.51 -> tools/gen1.py (auto); review=proofread-sonnet; flags=none
\begin{question}{40}
가등기담보 등에 관한 법률이 적용되는 양도담보의 효력에 관한 설명으로 옳지 않은 것은? (다툼이 있으면 판례에 의함)
\begin{choices}
  \item 양도담보권의 피담보채권의 범위는 저당권에 의해 담보되는 채권의 범위에 관한 규정이 적용된다.
  \item 양도담보권자에게는 양도담보권에 기한 물상대위가 인정된다.
  \item 양도담보권자가 담보목적 부동산을 임의로 처분한 경우, 그 부동산의 양수인이 선의이더라도 소유권을 취득하지 못한다.
  \item 양도담보권자가 귀속청산의 방법에 의하여 확정적으로 소유권을 취득하려면 적법하게 청산절차를 거쳐야 한다.
  \item 양도담보권은 우선변제적 효력이 있다.
\end{choices}
\end{question}
\answer{3}
